\documentclass{article}
\usepackage{graphicx}
\usepackage[margin=2cm, a4paper]{geometry}
\usepackage{setspace}
\usepackage{pdfpages}
\usepackage{float}
\usepackage{amsmath}
\usepackage{fancyvrb}
\usepackage{amssymb}
\usepackage{minted}
\usepackage{enumitem}
\usepackage[colorlinks,linkcolor=blue]{hyperref}


\renewcommand{\contentsname}{Contents}
\renewcommand{\figurename}{Figure}
\renewcommand{\tablename}{Table}
\hypersetup{
    colorlinks=true,
    linkcolor=black,
    filecolor=magenta,      
    urlcolor=blue,
}
\newcommand{\mytitle}{Abstract Algebra --- Homework 2}
\newcommand{\myauthor}{B13902022 Yu-Chi,Lai}

\usepackage{fancyhdr}
\pagestyle{fancy}
\fancyhead{}
\fancyhead[L]{\mytitle}
\fancyhead[R]{\myauthor}

\usepackage{coffeestains}

\title{\mytitle}
\author{\textbf{\myauthor}}
\date{}

\begin{document}

\maketitle

% \coffeestainB{0.6}{0.8}{30}{10cm}{15cm}
\section{}
\subsection*{(a)}
By the tower law, $[K:\mathbb{R}]=[K:\mathbb{C}][\mathbb{C}:\mathbb{R}]=2^{a+1}m$. Let $L$ be the Galois closure of $K$ over $\mathbb{R}$. Set $G=\mathrm{Gal}(L/\mathbb{R})$, so $|G|=[L:\mathbb{R}]$. It's obvious that $|G|=2^bm'$ for some $b\le a, m'|m, \text{m' is odd}$ (also by tower law).

Let $P\le G$ be a sylow 2-subgroup of $G$, so $[G:P]=m'$ is odd. Let $F$ be the corresponding intermediate field, we have $[F:\mathbb{R}]=[G:P]=m'$. Suppose $\alpha\in F\backslash \mathbb{R}$, and its minimal irreducible polynomial has degree $d$ ($d\ge 2$). Clearly, $d|[F:\mathbb{R}]$ so $d$ is odd. But any odd-degree polynomial has a real root, which condradicts the irreducibility. Hence, $F=\mathbb{R}$ and $G=P$ is a 2-group.

Therefore, $[L:\mathbb{R}]=2^b$, by the tower law, $[L:\mathbb{R}]=[L:\mathbb{C}][\mathbb{C}:\mathbb{R}]=2[L:\mathbb{C}]$, $[L:\mathbb{C}]=2^{b-1}$. Since $[K:\mathbb{C}]$ divides $[L:\mathbb{C}]$, the only possibility is $a=0$.
\subsection*{(b)}
Assume, $a\le 1$.
Now define $H=\mathrm{Gal}(L/\mathbb{C})\le G$, the subgroup is not trivial since $|H|=[L:\mathbb{C}]=[L:K][\mathbb{K}:\mathbb{C}]\le 2^a\le 2$

Since $2$ is the least prime dividing $|H|$, there exists a normal subgroup of $H$ of index $2$, let it be $N$, $[H:N]=2$. By the Galois correspondance, the fixed field $E=L^N$ satisfies $[E:\mathbb{C}]=[H:N]=2$, so $E$ is a degree-2 extension of $\mathbb{C}$.

Pick any $\beta\in E\backslash \mathbb{C}$, its minimal irreducible polynomial over $\mathbb{C}$ must be quadratic:
\[x^2+bx+c, b,c\in \mathbb{C}\]
By the formula, the roots are $\frac{-b\pm \sqrt{b^2-4ac}}{2}$, now $b^2-4ac\in \mathbb{C}$ and every complex number has its square root in $\mathbb{C}$ trivially, so $\beta\in \mathbb{C}$, a contradiction. 

Hence, $a=0$.
\section{}
\subsection*{(a)}
The minimal polynomial of $\sqrt[3]{2}$ over $\mathbb{Q}$ is $x^3-2=(x-\sqrt[3]{2})(x^2+2x+2)$, which does not split in $\mathbb{Q}(\sqrt[3]{2})$. Hence, $\mathbb{Q}(\sqrt[3]{2})/\mathbb{Q}$ is not a normal extension.
\subsection*{(b)}
Since $\mathbb{Q}(\sqrt[4]{2}, \zeta_8)/\mathbb{Q}$ contains all roots of $x^4-2$ ($\sqrt[4]{2}, -\sqrt[4]{2}, \zeta_8^2\sqrt[4]{2}, -\zeta_8^2\sqrt[4]{2}$), it's a splitting field for the irreducible polynomial (in $\mathbb{Q}[x]$) $x^4-2$. Every splitting field is normal, so $\mathbb{Q}(\sqrt[4]{2}, \zeta_8)/\mathbb{Q}$ is normal.

Clearly, $\mathbb{Q}(\sqrt[4]{2}, \zeta_8)=\mathbb{Q}(\sqrt[4]{2}, i)$. By tower law, $[\mathbb{Q}(\sqrt[4]{2}), i:\mathbb{Q}]=[\mathbb{Q}(\sqrt[4]{2}, i):\mathbb{Q}(i)][\mathbb{Q}(i):\mathbb{Q}]=4\times 2=8$. Hence, the degree for the extension is $8$.
\end{document}
% how to display codes?
% \begin{minted}[frame=lines,framesep=2mm,baselinestretch=1.2,linenos,breaklines]{python}
% \end{minted}

% how to display images?
% \begin{figure}[H]
%     \centering
%     \includegraphics[width=0.5\linewidth]{}
%     \caption{Caption}
% \end{figure}
% test